\section{Detection Algorithm and Per-Series Reasoner}
\label{sec:edge}

\begin{figure}[t]
  \centering
  \fitw{%
  \begin{tikzpicture}[node distance=5mm and 7mm]
    \node[io] (s) {sample};
    \node[box, right=of s] (cn) {counter?\\rate-normalize};
    \node[box, right=of cn] (wb) {rebuild baseline\\over window\\(Welford)};
    \node[box, right=of wb] (z) {$z=\dfrac{|x-\mu|}{\sigma}$};
    \node[dec, right=of z] (br) {$z\ge\tau$?};
    \node[box, above right=3mm and 9mm of br] (hold) {increment count;\\ anomalous if\\ count $\ge c$};
    \node[box, below right=3mm and 9mm of br] (reset) {reset count;\\ admit sample\\ to window};
    \node[io, right=22mm of br] (v) {verdict};
    \draw[ar] (s) -- (cn);
    \draw[ar] (cn) -- (wb);
    \draw[ar] (wb) -- (z);
    \draw[ar] (z) -- (br);
    \draw[ar] (br) |- node[near start,right,font=\scriptsize]{breach} (hold);
    \draw[ar] (br) |- node[near start,right,font=\scriptsize]{clean} (reset);
    \draw[ar] (hold) -| (v);
    \draw[ar] (reset) -| (v);
  \end{tikzpicture}}
  \caption{One detection pass. Counters are rate-normalized, the baseline is rebuilt
  over the window, the sample is scored, and the breach branch either confirms (and
  may emit \emph{anomalous}) or resets and admits the sample.}
  \label{fig:edge-pipeline}
\end{figure}

\subsection{Baseline and breach test}
For a series the reasoner keeps a window of the most recent clean samples, up to a
size $W$ (default $300$). Over the window it computes the mean $\mu$ and standard
deviation $\sigma$ with Welford's method~\citep{welford1962}. For an incoming
sample $x$ the breach score is the absolute $z$-score,
\begin{equation}
  z = \frac{\lvert x - \mu \rvert}{\sigma},
\end{equation}
and the sample breaches when $z \ge \tau$, with $\tau$ the threshold in standard
deviations (default $3$). The window must hold at least $m$ clean samples before
the reasoner will score against it (default $m = 30$); until then the series is
reported as \emph{insufficient baseline} rather than scored.

A flat baseline needs care, because $\sigma = 0$ makes the $z$-score undefined.
When the standard deviation is within floating-point noise of zero, the reasoner
returns $0$ if the sample equals the mean and otherwise a magnitude-aware score
that always clears $\tau$ and grows with the relative excursion, so any departure
from a perfectly flat series is flagged and larger departures rank above smaller
ones.

\subsection{Baseline integrity}
A breaching sample is not admitted to the window. Only clean samples update the
baseline. This single rule, taken from \dc{}'s DDoS example, keeps a sustained
anomaly from raising its own normal: a series that jumps and stays high keeps being
measured against the pre-jump baseline rather than slowly accepting the jump as
the new ordinary.

\subsection{Confirmation hysteresis}
A single breach is not yet an anomaly. The reasoner counts consecutive breaching
samples and only reports \emph{anomalous} once that run reaches a confirmation
length $c$ (default $5$). A breach before the run is long enough is reported as
\emph{pending anomaly}, and any clean sample resets the run to zero and is admitted
to the baseline. A pass therefore ends in one of four states: insufficient
baseline, clean, pending anomaly, or anomalous. The whole sequence (hydrate,
evaluate, branch on breach, finalize) is the \texttt{CausalFlow} pipeline of
\cref{sec:background}.

\begin{figure}[t]
  \centering
  \fitw{%
  \begin{tikzpicture}[node distance=15mm and 15mm]
    \node[sstate] (ib) {insufficient\\baseline};
    \node[sstate, right=of ib] (cl) {clean};
    \node[sstate, right=of cl] (pa) {pending\\anomaly};
    \node[sstate, right=of pa] (an) {anomalous};
    \draw[ar] (ib) -- node[above,font=\scriptsize]{$\ge m$ clean} (cl);
    \draw[ar] (cl) -- node[above,font=\scriptsize]{breach} (pa);
    \draw[ar] (pa) -- node[above,align=center,font=\scriptsize]{breach,\\count $\ge c$} (an);
    \draw[ar] (pa) to[bend left=20] node[below,font=\scriptsize]{clean} (cl);
    \draw[ar] (an) to[bend left=32] node[below,font=\scriptsize]{clean} (cl);
  \end{tikzpicture}}
  \caption{Verdict states. A series stays in \emph{insufficient baseline} until it
  has $m$ clean samples, then sits at \emph{clean}. A breach moves it to
  \emph{pending anomaly}, and once breaches reach the confirmation length $c$ it
  becomes \emph{anomalous}. Any clean sample returns it to \emph{clean} and is
  admitted to the baseline.}
  \label{fig:states}
\end{figure}

\subsection{Counter rate-normalization}
Cumulative monotonic counters are converted to a per-second rate before scoring, so
the detector sees a rate series rather than an ever-climbing total. The reasoner
keeps the last reading per series and derives the rate from the change in value
over the change in time between two readings on the same reset lineage. It emits no
rate on the first reading, on a reset (the lineage anchor changed), on
non-monotonic time, or across a gap longer than two hours, since a rate across any
of those is not meaningful. A decrease is salvaged as a 32-bit counter wrap when
the implied rate is plausible and dropped otherwise. The rate then feeds the same
$z$-score path as a gauge.

\subsection{What the reasoner persists}
The reasoner is stateless per call: it receives the window and returns a verdict
plus the next window. The per-series engine persists only the bounded window tail
and the confirmation counter, not a running accumulator, and rebuilds the mean and
variance from the window on each pass. Keeping just the window makes a series
trivial to checkpoint, and means the engine never has to trust an accumulator whose
stored summary might disagree with the samples it is paired with. The cost is that
each pass recomputes over the window, which is $O(W)$ rather than the $O(1)$ a
reversible accumulator would give. The shared core implements the reversible
accumulator too; \cref{sec:eval} measures what the recompute costs in practice.
